\textbf{原图本身是一棵树}

连接终端点集合 $S$ 的最小连通子图是唯一的，就是所有终端点之间路径的并。

\begin{itemize}
  \item 需要显式保留结构：按 DFS 序排序终端点，加入相邻点的 LCA 后建立虚树；虚树边 $(u,v)$ 的权值是原树距离 $d(u,v)$。
  \item 只求总边权：将终端点按 DFS 序排列为 $v_1,\ldots,v_k$，令 $v_{k+1}=v_1$，则
  \[
    W=\frac12\sum_{i=1}^{k}d(v_i,v_{i+1}).
  \]
  \item 只有三个终端点时，直接使用
  $W=(d(a,b)+d(b,c)+d(c,a))/2$。
\end{itemize}

原因是按 DFS 序绕终端点一圈时，最小连接子树中的每条边恰好被跨过两次。

\textbf{一般非负权图，终端点很少}

使用 Dreyfus--Wagner 子集 DP。
设终端点为 $t_0,\ldots,t_{k-1}$，定义
\[
  f[S][v]=\text{连接 }S\text{ 中终端点并在 }v\text{ 汇合的最小代价}.
\]

初始化 $f[2^i][t_i]=0$。
对每个状态 $S$ 依次进行两步：

\begin{enumerate}
  \item 在同一个汇合点合并两个子集：
  \[
    f[S][v]\gets\min_{A\subset S,\,A\ne\varnothing}
    \bigl(f[A][v]+f[S\setminus A][v]\bigr).
  \]
  \item 把所有 $f[S][v]$ 作为多源初值跑一次 Dijkstra，使汇合点可以沿最短路移动。
\end{enumerate}

答案是 $\min_v f[2^k-1][v]$。
若需要输出方案，为“子集合并”和“最短路转移”分别记录前驱。
一般图斯坦纳树是 NP-hard，这个算法只适合 $k$ 很小；状态解释可参考
\href{https://codeforces.com/blog/entry/91601}{Codeforces：Dreyfus--Wagner DP}。
